\begin{figure}[htbp]
\centering
\begin{tikzpicture}[x=1.7cm,y=0.55cm,font=\scriptsize]
\foreach \k/\lab in {0/{PUSH 6},1/{PUSH 7},2/{MUL},3/{PUSH 8},4/{ADD}}{
  \node[anchor=south] at (\k+0.5,2.4) {\lab};}
\foreach \k/\a in {0/6,1/6,2/42,3/42,4/50}{
  \draw[fill=accentsoft,draw=accent] (\k+0.1,0) rectangle (\k+0.9,1);
  \node at (\k+0.5,0.5) {\a};}
\foreach \k/\b in {1/7,3/8}{
  \draw[fill=goodbg,draw=good] (\k+0.1,1) rectangle (\k+0.9,2);
  \node at (\k+0.5,1.5) {\b};}
\node[anchor=east] at (0.05,0.5) {bottom};
\node[anchor=east] at (0.05,1.5) {next};
\node[anchor=north,text=slate] at (2.5,-0.2) {MUL pops 7 and 6, pushes 42; ADD pops 8 and 42, pushes 50};
\end{tikzpicture}
\caption{The operand stack of \texttt{ch18\_vmrun} after each step of
$6\times7+8$ (the real trace is in the worked example). Operands are
never named: the top of the stack is always ``the current value''.}
\label{fig:trace18}
\end{figure}
